\subsection{绝对半单代数的上同调}

命题 6. —— 设 K 是一个交换环，A 是一个 K-代数。设 \(e = \sum_{i=1}^{r} a_i \otimes a_i'\) 是 \(B = A \otimes_K A^o\) 的一个满足 VIII, 第 236 页的命题 5 中 (ii) 的条件的元素。对任意整数 \(n \geqslant 1\) 与 \(C^n(A, P)\) 的任意元素 f，我们用 \(\gamma^n(f)\) 表示由下述公式定义的 \(C^{n-1}(A, P)\) 的元素

\[
\gamma^ {n} (f) \left(x _ {1}, \dots , x _ {n - 1}\right) = \sum_ {i = 1} ^ {r} a _ {i} f \left(a _ {i} ^ {\prime}, x _ {1} \dots , x _ {n - 1}\right).\tag{15}
\]

于是我们有

\[
\partial^ {n - 1} (\gamma^ {n} (f)) + \gamma^ {n + 1} (\partial^ {n} (f)) = f\tag{16}
\]

对每个整数 \(n \geqslant 1\) 与每个 \(f \in \mathrm{C}^n(\mathrm{A},\mathrm{P})\) 成立。

*注 1. —— 态射 \(\partial^n: \mathrm{C}^n(\mathrm{A},\mathrm{P}) \to \mathrm{C}^{n+1}(\mathrm{A},\mathrm{P})\) 定义了 K-模的一个复形 \((\mathrm{C}(\mathrm{A},\mathrm{P}),\partial)\)（X, §2, n° 1, 第 24 页）。于是映射 \(\gamma_n\) 定义了从该复形到它自身、把 0 连接到 \(\mathrm{Id}_{\mathrm{C}(\mathrm{A},\mathrm{P})}\) 的一个同伦（X, §2, n° 4, 第 32 页，定义 4）。 *

我们沿用 No.6 的记号。对每个整数 \(n \geqslant 0\)，我们用公式定义映射 \(h_n: \mathrm{B}_n \to \mathrm{B}_{n+1}\)

\[
h _ {n} (x) = d _ {n + 2} ^ {1} (e \otimes x) = \sum_ {i = 1} ^ {r} a _ {i} \otimes a _ {i} ^ {\prime} x.
\]

它是 \((A, A)\) -双模的一个同态（公式 (3)）。

引理 5. —— 我们有关系式

\[
d _ {n + 1} \circ h _ {n} + h _ {n - 1} \circ d _ {n} = 1 _ {\mathrm{B} _ {n}}\tag{17}
\]

对每个 \(n \geqslant 1\) 成立。

设 \(x\in \mathrm{B}_n\)；我们有

\[
\begin{array}{c} (d _ {n + 1} \circ h _ {n}) (x) = (d _ {n + 1} \circ d _ {n + 2} ^ {1}) (e \otimes x) \\ = (d _ {n + 1} ^ {0} \circ d _ {n + 2} ^ {1}) (e \otimes x) - \sum_ {i = 2} ^ {n + 2} (- 1) ^ {i} (d _ {n + 1} ^ {i - 1} \circ d _ {n + 2} ^ {1}) (e \otimes x), \end{array}
\]

于是，由公式 (6)，

\[
(d _ {n + 1} \circ h _ {n}) (x) = (d _ {n + 1} ^ {0} \circ d _ {n + 2} ^ {0}) (e \otimes x) - \sum_ {i = 2} ^ {n + 2} (- 1) ^ {i} (d _ {n + 1} ^ {1} \circ d _ {n + 2} ^ {i}) (e \otimes x).
\]

但我们有

\[
(d _ {n + 1} ^ {0} \circ d _ {n + 2} ^ {0}) (e \otimes x) = \varepsilon (e) x = x
\]

这由 VIII, 第 236 页的命题 5 的性质 (ii) 得出，而对 \(i \geqslant 2\)，

\[
d _ {n + 2} ^ {i} (e \otimes x) = e \otimes d _ {n} ^ {i - 2} (x),
\]

由此给出

\[
(d _ {n + 1} \circ h _ {n}) (x) = x - d _ {n + 1} ^ {1} (e \otimes d _ {n} (x)) = x - h _ {n - 1} \circ d _ {n} (x)
\]

从而得到公式 (17)。

利用引理 5，我们可以完成命题 6 的证明。设 n 是一个整数 \(\geqslant 1\)，f 是 \(C^{n}(A,P)\) 的一个元素。按构造，我们有

\[
\alpha^ {n - 1} (\gamma^ {n} (f)) = \alpha^ {n} (f) \circ h _ {n - 1}\tag{18}
\]

于是，由公式 (9) 与 (18)，

\[
\begin{array}{r l} \alpha^ {n} (\partial^ {n - 1} (\gamma^ {n} (f)) + \gamma^ {n + 1} (\partial^ {n} (f))) & = \alpha^ {n - 1} (\gamma^ {n} (f)) \circ d _ {n} + \alpha^ {n + 1} (\partial^ {n} (f)) \circ h _ {n} \\ & = \alpha^ {n} (f) \circ h _ {n - 1} \circ d _ {n} + \alpha^ {n} (f) \circ d _ {n + 1} \circ h _ {n} \\ & = \alpha^ {n} (f), \end{array}
\]

其中最后的等式由 (17) 得出。由于 \(\alpha^{n}\) 是双射，命题得证。

定理 3. —— 设 K 是一个交换环，A 是一个 K-代数，P 是一个 (A,A)-双模。假设 \((\mathrm{A}\otimes_{\mathrm{K}}\mathrm{A}^{\circ})\) -模 A 是投射的。则我们有 \(\mathrm{H}^{n}(\mathrm{A},\mathrm{P})=0\)（对每个整数 \(n\geqslant1\)）。

我们必须证明，对每个整数 \(n \geqslant 1\)，\(C^{n}(A, P)\) 的每个满足 \(\partial^{n}(f) = 0\) 的元素 f 都具有形式 \(\partial^{n-1}(g)\)（对 \(C^{n-1}(A, P)\) 的某个元素 g）。由命题 5，这是命题 6 的一个直接推论。

推论. —— 从 A 到 P 的每个 K-导子都是内导子。

这是等式 \(H^{1}(A,P)=0\) 的一个转述。

注. —— 2) 定理 3 的假设特别在 K 是一个域且 A 是一个绝对半单 K-代数时被满足（VIII, 第 238 页，定理 2）。

*3) 假设 K-模 A 是投射的。定理 3 也可按下述方式证明。复形 \((\bigoplus_{n\geqslant 0}\mathrm{B}_n,d)\) 与同态 \(\varepsilon :\mathrm{B}_0\to \mathrm{A}\) 定义了 B-模 A 的一个投射分解；因此，对每个 \(n\geqslant 0\)，K-模 \(\mathrm{H}^n (\mathrm{A},\mathrm{P})\) 同构于 \(\operatorname {Ext}_{\mathrm{B}}^{n}(\mathrm{A},\mathrm{P})\)（X, §6, n° 1, 第 100 页，定理 1）。若 B-模 A 是投射的，则 K-模 \(\operatorname {Ext}_{\mathrm{B}}^{n}(\mathrm{A},\mathrm{P})\) 对 \(n\geqslant 1\) 为零（X, §5, n° 3, 第 88 页，命题 5 的推论），由此推出 \(\mathrm{H}^n (\mathrm{A},\mathrm{P})\) 为零。反之，若对每个 (A, A)-双模 P 都有 \(\mathrm{H}^1 (\mathrm{A},\mathrm{P})\) 为零，则 B-模 A 是投射的（X, §5, n° 5, 第 93 页，命题 10）。*

